\section{Human authorship, integrated coverage and naming policy}
\label{sec:editorial-front}
\textbf{Human author: Hruday N M (BUNZEEY).} This is an integrated manuscript prepared with AI assistance in research, derivation checking, code, source comparison, and editing. Model-agent reviews are disclosed in the historical chapters. No external human peer review or formal proof-assistant verification is claimed; there is no asserted arXiv submission or acceptance.

The complete Draft01 research body and appendices through Round26 and the complete Round27 and Round28 research addenda are carried forward here, including failures, conditions and retrospective roadmaps. The original source packages remain unchanged. A dated roadmap is a record of the question selected then, not a claim that subsequently completed work remains unexecuted. The Round29 chapter retains its reviewed investigations and dated roadmap. The Round30 chapter contains the current three-investigation cycle and its closing roadmap. Equation numbers and cross-references are unified; original identifiers remain recoverable in the registry.

\subsection{What Hruday and HNM do and do not name}
Hruday names and the HNM namespace identify this project's records. They do not rename Yang--Mills theory, Haar measure, Wilson observables, Jensen's inequality, the GNS construction, Schur complements, Duhamel's formula or any established external theorem. Mathematical notation and physical units stay intact. In particular, adding HNM to a result does not convert a conditional estimate into a theorem about another model, and does not establish historical priority. The number of newly asserted axioms is zero. HNM-T and HNM-P statement aliases identify admitted scoped theorem statements and propositions by mathematical content, including source-theorem applications. These aliases add no premises or proof steps. External axioms retain their names and attribution.

Every contribution has a stable HNM-C identifier with its legacy ID. Labelled equations have stable HNM-E keys derived from their original source labels; the printed HNM equation numbers are document locators. HNM-Q quantity entries give traceable names while retaining their conventional symbols. Shared symbols with different meanings remain local to the stated model. The exact machine-readable registry is \repo{registry/hnm-registry.json}; the appendix prints the contribution aliases and equation concordance.

The double-star marker continues to mean derived, specialized, checked or implemented within this research record. It is not evidence that a method or equation is new to the literature. The priority table ranks usefulness to this route qualitatively, rather than assigning discovery scores or a percentage of the Millennium problem.

\subsection{The requested primary-source comparison changes the novelty boundary}
The September 2026 version of Gauvin's preprint \cite{gauvin2026} gives a regulated $\mathrm{SU}(3)$ analysis in Sections II.3--II.4, with supplemental proofs: physical gap estimates, local Wilson survival, a finite energy moment, and physical GNS correlation bounds. This overlaps the general strategy of our AJ/AK branch. The Wilson variance--moment--window--Jensen mechanism therefore must not be presented as a Hruday invention. Our stated $\mathrm{SU}(2)$ model, constants, domains, boundary convention and source-bound computations remain individually inspectable; their distinct scientific novelty is unverified. The preprint's constants cannot be transferred by relabeling its group or energy scales. Its abstract also leaves the interacting continuum construction open. This is a scoped comparison, not an endorsement or complete proof audit of the external work.

The numerical-cap AQ branch has a further explicit prior-method boundary: Gauvin's Supplement A.10 already presents the thermodynamic compactness, dynamics, GNS and Fourier-test strategy. The actual unbounded-onsite dynamics placement uses Nachtergaele--Sims~\citep{nachtergaele2014dynamics}, with the source panel also checking the modern Nachtergaele--Sims--Young framework~\citep{nachtergaele2019quasilocality}. The contribution is a separately checked $\mathrm{SU}(2)$ selected-strip model dictionary and numerical-cap application. Neither renaming nor a bounded-spin theorem supplies the unbounded-rotor hypotheses. Volume-norm convergence of a fixed local evolution remains distinct from norm continuity in time on every bounded local operator.

The official problem description \cite{jaffe_witten} remains the governing target: a nontrivial four-dimensional quantum theory of the requisite axiomatic strength with a positive physical mass gap for each compact simple gauge group. Finite graphs, a fixed-spacing homogeneous limit, a selected boundary state and a conditional spectral estimate do not jointly supply that construction. The pending model map, scale trajectory, uniform estimates, nontrivial observables and limiting reconstruction are explicit proof obligations.

\subsection{How to read the enlarged manuscript}
The historical first part establishes the model vocabulary and source-bound calculations. Round27 studies identifiability and one nonlinear correction. Round28 restores complete sources, improves finite-graph heat control and establishes conditional physical Wilson estimates. Round29 tests the next missing implications against the reviewed external sources. Use the ranked appendix to locate a useful result, then read its full assumptions and surviving objection. A negative or insufficient gate is a completed investigation with a limited conclusion, not a solved target.

\subsection{Current extension to the numerical-cap state}
The later AT1 gate closes the Wilson operator-domain and first-moment identity obligations directly in the AQ numerical-cap state, with second moment bounded by $\alpha^2(36+98|\tau|)$. This is a fresh same-state proof, using seven bulk stars, rather than a transfer from the earlier orthant state. Earlier passages saying that AQ has no moment theorem describe their original admissions; the current registry attaches the separate AT1 extension. The full forward and reverse derivations appear in the Round30 part. AO/AQ identification, second-moment equality, boundary-independent uniqueness and the continuum construction remain open. This distinction applies throughout the carried historical chapters and appendices.
